
\documentclass{article}
\usepackage{colortbl}
\usepackage{makecell}
\usepackage{multirow}
\usepackage{supertabular}

\begin{document}

\newcounter{utterance}

\centering \large Interaction Transcript for game `taboo', experiment `medium\_en', episode 3 with reflexion{-}Qwen3.5{-}4B.
\vspace{24pt}

{ \footnotesize  \setcounter{utterance}{1}
\setlength{\tabcolsep}{0pt}
\begin{supertabular}{c@{$\;$}|p{.15\linewidth}@{}p{.15\linewidth}p{.15\linewidth}p{.15\linewidth}p{.15\linewidth}p{.15\linewidth}}
    \# & $\;$A & \multicolumn{4}{c}{Game Master} & $\;\:$B\\
    \hline

    \theutterance \stepcounter{utterance}  
    & & \multicolumn{4}{p{0.6\linewidth}}{
        \cellcolor[rgb]{0.9,0.9,0.9}{
            \makecell[{{p{\linewidth}}}]{
                \texttt{\tiny{[P1$\langle$GM]}}
                \texttt{You have attempted this same task before and did not succeed. These are the plans you wrote after those attempts.} \\
\texttt{Trial \#1: New plan: I failed because I provided clues that were too generic and allowed the opponent to guess related words (Offense, Misdemeanor) instead of the target word (violation). I need to shift my strategy from defining the concept broadly to using specific, unique identifiers that distinguish "violation" from its synonyms like "transgression" or "offence." Specifically, I should focus on the legal context of "breaking a law" without using the word "law" itself, or emphasize the specific consequence of "penalty" in a way that excludes other related terms. I will avoid defining the crime by its severity (less serious than felony) as that led to "misdemeanor." Instead, I will use a clue that highlights the specific nature of the act or a unique legal term associated only with "violation" to force the opponent to guess the exact word within 3 tries.} \\
\\ 
\texttt{Trial \#2: New plan: I failed because my clues were too broad and focused on severity or general legal concepts, allowing the opponent to guess related words like "misdemeanor" or "infraction." I need to shift my strategy to use a specific, unique identifier that only applies to "violation" and excludes all related terms. Specifically, I should focus on the concept of a "technical" or "minor" breach that is often cited on a **driver's license** or **DMV record**, as this is a unique context for "violation" that does not apply to "transgression" or "offence." I will avoid mentioning "law," "crime," "fine," or "penalty" in a way that suggests other categories, and instead emphasize the administrative nature of the record.} \\
\\ 
\texttt{Trial \#3: New plan: I failed because my clues repeatedly triggered the opponent to guess related words like "offense" or "misdemeanor" by focusing on general legal concepts or severity. I need to adopt a strategy that isolates "violation" through its specific administrative consequence: the immediate addition of points to a **driver's license** or **DMV record**. I will construct a clue that explicitly mentions "points on a driving record" and "minor traffic breach" while strictly avoiding synonyms like "offense," "infraction," "breach," or "law." The goal is to make the description so specific to the context of a traffic ticket on a license that "violation" is the only logical answer, preventing the opponent from guessing broader legal terms.} \\
\\ 
\\ 
\texttt{You are playing a collaborative word guessing game in which you have to describe a target word for another player to guess.} \\
\\ 
\texttt{Rules:} \\
\texttt{(a) You have to reply in the form: CLUE: <some text>. Guesses from the other player will start with GUESS.} \\
\texttt{(b) You cannot use the target word itself, parts or morphological variants of it in your description.} \\
\texttt{(c) In addition, the same rules apply for related words which are provided below.} \\
\\ 
\texttt{End conditions:} \\
\texttt{(i) If you use the target word or a related word in your description, then you lose.} \\
\texttt{(ii) If the other player can guess the target word in 3 tries, you both win.} \\
\\ 
\texttt{Let us start.} \\
\\ 
\texttt{This is the target word that you need to describe and that the other player needs to guess:} \\
\\ 
\texttt{violation} \\
\\ 
\texttt{Related words are:} \\
\\ 
\texttt{{-} transgression} \\
\texttt{{-} misdemeanor} \\
\texttt{{-} offence} \\
\\ 
\texttt{Important: You are under time pressure, give short descriptions that are to the point!} \\
            }
        }
    }
    & \\ \\

    \theutterance \stepcounter{utterance}  
    & \multicolumn{4}{p{0.6\linewidth}}{
        \cellcolor[rgb]{0.9,0.9,0.9}{
            \makecell[{{p{\linewidth}}}]{
                \texttt{\tiny{[P1$\rangle$GM]}}
                \texttt{CLUE: A minor traffic breach that results in points being added to a driver's license.} \\
            }
        }
    }
    & & \\ \\

    \theutterance \stepcounter{utterance}  
    & & & \multicolumn{2}{p{0.3\linewidth}}{
        \cellcolor[rgb]{0.9,0.9,0.9}{
            \makecell[{{p{\linewidth}}}]{
                \texttt{\tiny{[GM$|$GM]}}
                \texttt{continue} \\
            }
        }
    }
    & & \\ \\

    \theutterance \stepcounter{utterance}  
    & & & \multicolumn{2}{p{0.3\linewidth}}{
        \cellcolor[rgb]{0.9,0.9,0.9}{
            \makecell[{{p{\linewidth}}}]{
                \texttt{\tiny{[GM$|$GM]}}
                \texttt{a minor traffic breach that results in points being added to a drivers license} \\
            }
        }
    }
    & & \\ \\

    \theutterance \stepcounter{utterance}  
    & & \multicolumn{4}{p{0.6\linewidth}}{
        \cellcolor[rgb]{0.9,0.9,0.9}{
            \makecell[{{p{\linewidth}}}]{
                \texttt{\tiny{[GM$\rangle$P2]}}
                \texttt{You have attempted this same task before and did not succeed. These are the plans you wrote after those attempts.} \\
\texttt{Trial \#1: New plan: I failed because I treated the clues as definitions of the target word itself rather than as descriptors of its category. The first clue described the general concept ("an act that breaks a law"), which I answered with "Offense." The second clue narrowed the scope to a specific type of crime ("less serious than a felony"), which I answered with "Misdemeanor." However, the third clue ("You can be charged with this if you break a law, but it is not considered a serious felony") is logically identical to the second clue, meaning the target word must be the specific category name that encompasses both "Misdemeanor" and "Offense" as examples, not a specific instance of the crime. I should have recognized that the clues were defining the class of words (e.g., "Crime," "Violation," "Transgression") rather than guessing the specific word directly. My next guess should be a general term for illegal acts, such as "Crime."} \\
\\ 
\texttt{Trial \#2: New plan: I failed because I treated the clues as direct definitions of the target word itself, guessing specific types of crimes (Misdemeanor, Infraction) instead of the broader category that encompasses them. The clues describe the nature of the act (a minor legal infraction, a breach involving fines but not imprisonment) rather than naming a specific crime. I should have recognized that the target word is the general class of these actions, such as "Violation" or "Offense," rather than a specific instance within that class. My next guess should be a general term for illegal acts, such as "Violation."} \\
\\ 
\texttt{Trial \#3: New plan: I failed because I guessed specific types of traffic violations (like "Ticket") instead of the general category name that encompasses all such infractions. The clues describe the nature of the act (a minor legal infraction, a breach involving fines) rather than naming a specific instance. I should have recognized that the target word is the broad class of these actions, such as "Violation" or "Infraction," rather than a specific document or type of crime. My next guess should be the general term for minor traffic offenses, such as "Infraction."} \\
\\ 
\\ 
\texttt{You are playing a collaborative word guessing game in which you have to guess a target word that another player describes to you.} \\
\\ 
\texttt{You can make one guess at each trial. You win when you guess the target word. You lose when you cannot guess it in 3 tries.} \\
\\ 
\texttt{After each trial you will get a new hint from the other player which starts with CLUE.} \\
\\ 
\texttt{Make your guesses by just saying the word using the following form: GUESS: <a word>} \\
\\ 
\texttt{Let us start.} \\
\\ 
\\ 
\texttt{CLUE: A minor traffic breach that results in points being added to a driver's license.} \\
            }
        }
    }
    & \\ \\

    \theutterance \stepcounter{utterance}  
    & & & \multicolumn{4}{p{0.6\linewidth}}{
        \cellcolor[rgb]{0.9,0.9,0.9}{
            \makecell[{{p{\linewidth}}}]{
                \texttt{\tiny{[GM$\langle$P2]}}
                \texttt{GUESS: Infraction} \\
            }
        }
    }
     \\ \\

    \theutterance \stepcounter{utterance}  
    & & & \multicolumn{2}{p{0.3\linewidth}}{
        \cellcolor[rgb]{0.9,0.9,0.9}{
            \makecell[{{p{\linewidth}}}]{
                \texttt{\tiny{[GM$|$GM]}}
                \texttt{continue} \\
            }
        }
    }
    & & \\ \\

    \theutterance \stepcounter{utterance}  
    & & & \multicolumn{2}{p{0.3\linewidth}}{
        \cellcolor[rgb]{0.9,0.9,0.9}{
            \makecell[{{p{\linewidth}}}]{
                \texttt{\tiny{[GM$|$GM]}}
                \texttt{infraction} \\
            }
        }
    }
    & & \\ \\

    \theutterance \stepcounter{utterance}  
    & & \multicolumn{4}{p{0.6\linewidth}}{
        \cellcolor[rgb]{0.9,0.9,0.9}{
            \makecell[{{p{\linewidth}}}]{
                \texttt{\tiny{[P1$\langle$GM]}}
                \texttt{GUESS: Infraction} \\
            }
        }
    }
    & \\ \\

    \theutterance \stepcounter{utterance}  
    & \multicolumn{4}{p{0.6\linewidth}}{
        \cellcolor[rgb]{0.9,0.9,0.9}{
            \makecell[{{p{\linewidth}}}]{
                \texttt{\tiny{[P1$\rangle$GM]}}
                \texttt{CLUE: It is a specific type of minor traffic offense that gets recorded on a DMV, distinct from a misdemeanor or a felony.} \\
            }
        }
    }
    & & \\ \\

    \theutterance \stepcounter{utterance}  
    & & & \multicolumn{2}{p{0.3\linewidth}}{
        \cellcolor[rgb]{0.9,0.9,0.9}{
            \makecell[{{p{\linewidth}}}]{
                \texttt{\tiny{[GM$|$GM]}}
                \texttt{continue} \\
            }
        }
    }
    & & \\ \\

    \theutterance \stepcounter{utterance}  
    & & & \multicolumn{2}{p{0.3\linewidth}}{
        \cellcolor[rgb]{0.9,0.9,0.9}{
            \makecell[{{p{\linewidth}}}]{
                \texttt{\tiny{[GM$|$GM]}}
                \texttt{Related word 'misdemeanor' (stem=misdemeanor) is similar to clue word 'misdemeanor' (stem=misdemeanor)} \\
            }
        }
    }
    & & \\ \\

    \theutterance \stepcounter{utterance}  
    & \multicolumn{4}{p{0.6\linewidth}}{
        \cellcolor[rgb]{0.95,0.95,0.95}{
            \makecell[{{p{\linewidth}}}]{
                \texttt{\tiny{[P2$\rightarrow$P1]}}
                \texttt{Post{-}episode Reflexion (not sent to GM):} \\
\texttt{New plan: I failed because my previous clues used synonyms like "breach," "offense," and "minor traffic" which allowed the opponent to guess related words like "infraction" or "misdemeanor." I need to adopt a strategy that relies exclusively on the unique administrative mechanism of the target word: the specific action of **adding points** to a **driver's license**. I will construct a single, highly specific clue that mentions "points on a driver's license" and "traffic record" while strictly avoiding any words that could trigger synonyms like "fine," "penalty," "offense," "breach," or "law." The goal is to make the description so uniquely tied to the DMV point system that "violation" is the only possible answer, preventing the opponent from guessing broader legal terms.} \\
            }
        }
    }
    & & \\ \\

    \theutterance \stepcounter{utterance}  
    & & & \multicolumn{4}{p{0.6\linewidth}}{
        \cellcolor[rgb]{0.95,0.95,0.95}{
            \makecell[{{p{\linewidth}}}]{
                \texttt{\tiny{[P1$\rightarrow$P2]}}
                \texttt{Post{-}episode Reflexion (not sent to GM):} \\
\texttt{New plan: I failed because I treated the clues as definitions of the target word itself rather than as descriptors of its category. The clues describe the general concept ("a minor traffic breach"), which I answered with "Infraction." However, the clues are logically identical to the definition of the broader category, meaning the target word must be the specific category name that encompasses "Infraction" as an example, not a specific instance of the crime. I should have recognized that the clues were defining the class of words (e.g., "Traffic Violation," "Moving Violation," "Offense") rather than guessing the specific word directly. My next guess should be a general term for illegal traffic acts, such as "Violation."} \\
            }
        }
    }
     \\ \\

\end{supertabular}
}

\end{document}
